\subsection{紧李代数}

命题 1. 令 \(\mathfrak{g}\) 为（实）李代数。下列条件等价：

\begin{enumerate}
\def\labelenumi{(\roman{enumi})}
\item
  g 同构于某个紧李群的李代数。
\item
  群 \(\operatorname{Int}(\mathfrak{g})\)（第III章§6第2节定义2）是紧的。
\item
  \(\mathfrak{g}\) 有一个对称、正且分离的不变双线性形式（第I章§3第6节）。
\item
  g 是约化的（第I章§6第4节定义4）；对所有 \(x \in g\)，自同态 ad x 半单，且特征值都是纯虚数。
\item
  \(\mathfrak{g}\) 约化且其 Killing 型 B 是负的。
\item
  \(\Longrightarrow\) (ii)：若 g 是紧李群 G 的李代数，则群 \(\operatorname{Int}(\mathfrak{g})\) 是分离的，且同构于紧群 \(G_{0}\) 的一个商（第III章§6第4节推论4），因而是紧的。
\item
  \(\Longrightarrow\) (iii)：若群 \(\operatorname{Int}(\mathfrak{g})\) 紧，则 g 上存在一个对称双线性形式，它正、分离且在 \(\operatorname{Int}(\mathfrak{g})\) 下不变（第1节），从而也在 g 的伴随表示下不变。
\item
  \(\Longrightarrow\) (iv)：若 (iii) 成立，则 g 的伴随表示半单（第1节），因而 g 约化；此外，\(x \in g\) 的自同态 ad x 具有所述性质（第1节）。
\item
  \(\Longrightarrow\) (v)：对所有 \(x\in \mathfrak{g}\)，\(\mathrm{B}(x,x) = \mathrm{Tr}((\mathrm{ad}x)^2)\)；于是 \(\mathrm{B}(x,x)\) 是 ad \(x\) 的特征值的平方和，因而当这些特征值都是纯虚数时它是负的。
\item
  \(\Longrightarrow\) (i)：假设 g 约化，从而是一个交换子代数 c 与一个半单子代数 s 的积（第I章§6第4节命题5）。s 的 Killing 型是形式 B 在 s 上的限制，因而当 B 负时它是负且分离的。子群 \(\operatorname{Int}(\mathfrak{s})\)（\(\mathbf{GL}(\mathfrak{s})\) 的）是闭的（它是 \(\operatorname{Aut}(\mathfrak{s})\) 的单位元连通分支，第III章§10第2节推论2）且使分离的正形式 -B 不变；于是它是紧的，且 s 同构于紧李群 \(\operatorname{Int}(\mathfrak{s})\) 的李代数。再者，由于 c 交换，它同构于某个环面 T 的李代数。于是 g 同构于紧李群 \(\operatorname{Int}(\mathfrak{s}) \times \operatorname{T}\) 的李代数。
\end{enumerate}

定义 1. 紧李代数\(^{4}\)是指具有命题 1 中性质 (i) 至 (v) 的李代数。

于是，紧李代数就是交换代数与紧半单代数的积。换句话说，李代数紧当且仅当它约化且其导李代数紧。

紧李群的李代数是紧的。

命题 2. a) 有限个李代数的积是紧李代数，当且仅当每个因子都紧。

\begin{enumerate}
\def\labelenumi{\alph{enumi})}
\setcounter{enumi}{1}
\item
  紧李代数的子代数是紧的。
\item
  令 \(\mathfrak{h}\) 为紧李代数 \(\mathfrak{g}\) 的理想。则代数 \(\mathfrak{g} / \mathfrak{h}\) 紧，且扩张 \(\mathfrak{h} \to \mathfrak{g} \to \mathfrak{g} / \mathfrak{h}\) 平凡。
\end{enumerate}

断言 \(a\) 与 \(b\) 由命题 1 的刻画 (iii) 得出。部分 \(c\) 由 \(a\) 以及在约化李代数中每个理想都是直积因子这一事实（第I章§6第4节命题5的推论）得出。

命题 3. 令 G 为连通分支群有限的李群。下列条件等价：

\begin{enumerate}
\def\labelenumi{(\roman{enumi})}
\item
  李代数 \(\mathrm{L}(\mathrm{G})\) 紧。
\item
  群 \(\mathrm{Ad}(\mathrm{G})\) 紧。
\item
  L(G) 上存在在 G 的伴随表示下不变的分离的正对称双线性形式。
\end{enumerate}

*（iv）G 有在左平移与右平移下不变的黎曼度量。*

\begin{enumerate}
\def\labelenumi{(\roman{enumi})}
\item
  \(\Longrightarrow\) (ii)：若 L(G) 紧，则群 \(\mathrm{Ad}(\mathrm{G}_{0}) = \mathrm{Int}(\mathrm{L}(\mathrm{G}))\) 紧；由于它在 \(\mathrm{Ad}(\mathrm{G})\) 中的指标有限，后一群也紧。
\item
  \(\Longrightarrow\) (iii)：这由第1节得出。
\item
  \(\Longrightarrow\) (i)：由于 \(\operatorname{Int}(\mathrm{L}(\mathrm{G})) \subset \operatorname{Ad}(\mathrm{G})\)，这由命题 1 的刻画 (iii) 得出。
\end{enumerate}

\(^{*}\) (iii) \(\Longleftrightarrow\) (iv)：这由第III章§3第13节得出。*

